\section{Christmas and the Winter Solstice}

In `Roma e il natale solare nella tradizione nordico-aria' (\emph{La Difesa della razza}, 1940), Evola writes:

\begin{quotationx}
Very few suspect that the holidays i.e., Catholic holy days of today, in the century of skyscrapers, radio, great movements of the masses, are celebrated and continue ... a remote tradition, bringing us back to the times when, almost at the dawn of humanity, the rising motion of the first Aryan civilisation began; a tradition, in which, moreover, the great voice of those men is expressed rather than a particular belief.

A fact unknown to most must be first of all remembered, viz., that in its origins the date of \textbf{Christmas} and that of the beginning of the new year coincided, this date not being arbitrary, but connected to a precise cosmic event, namely, the \textbf{Winter Solstice}. The Winter Solstice falls in fact on December 25, which is the date of Christmas, subsequently known, but which in its origins had an essentially solar significance. That appears also in ancient Rome: the date of Christmas in ancient Rome was that of the rising of the Sun, the unconquered God, \emph{Natalis solis invicti}. With that, as the day of the new sun --- \emph{dies solis} novi --- in the imperial epoch brought the beginning of the new year, the new cycle. But this ``solar birth'' of Rome in the imperial period, in its turn, referred to a somewhat more remote tradition of Nordic-Aryan origin. Of the reset, Sol, the solar divinity, appeared already among the dii indigetes, that is, among the divinities of Roman origin, passed on from even more distant cycles of civilisation. In reality, the solar religion of the imperial period, in a large measure had the meaning of a recovery and almost of a rebirth, unfortunately altered by various factors of decomposition, of a very old Aryan heritage.
\end{quotationx}

First of all, let us clear up a misconception of the date December 25, which some have believed represents an error on Evola's part. \textbf{Julius Caesar}, in his calendar reform of 46 BC, did indeed set the date for the Winter Solstice on December 25. Over the centuries, the subtle error in the length of the Julian year gradually moved the astronomical date backwards. When \textbf{Pope Gregory} reformed the Julian calendar, he began with a later year, corresponding to a Church council, rather than the original year of 46 BC. Thus, the Winter Solstice moved back to around December 21 instead of the original December 25. Thus, Evola was correct in his claim that the data of Christmas was set to coincide with the Winter Solstice.

In the article, Evola goes on to tie in the date with the birth of Mithras. He also points out that in ancient Rome, the ``day of the sun'' was also the ``Lord's day'', another tradition adopted by the Christians. He also ties this symbolism of the light with the prologue to John's Gospel.

\begin{quotex}
The true light that gives light to every man was coming into the world
\end{quotex}


Evola concludes the article

\begin{quotex}
In the Aryan and Nordic tradition and in Rome itself, the same theme had an importance not only religious and mystical, but sacred, heroic and cosmic at the same time. It was the tradition of a people, to whom the same nature, the same great voice, which I wrote about , at that date, a tradition of a mystery of resurrection, of the birth or rebirth of a beginning not only of ``light'' and new life, but also of Imperium, in the highest and most august meaning of the word.
\end{quotex}



\flrightit{Posted on 2010-12-28 by Cologero}

\begin{center}* * *\end{center}

\begin{footnotesize}
\begin{sffamily}
\texttt{Charlotte Cowell on 2011-12-20 at 13:22 said:}

Magicians are supposed to rest on the winter solstice :-)

\hfill

\texttt{Cologero on 2011-12-20 at 14:53 said:}

Thanks my plan, Charlotte.

\hfill

\texttt{Charlotte Cowell on 2011-12-21 at 04:33 said:}

Happy Hannukah! May the spiritual lights burn brighter than ever in these days of darkness. Cxx

\hfill

\texttt{Attila on 2011-12-22 at 23:31 said:}

I was informed by a neighbor that last night was celebrated as Shab-e Yalda in Iran- which ostensibly represents the turning point -where darkness starts to give way to light- probably in connection with the birth of Mithra- the Sun god of ancient Persia. It's a beautiful holiday in the line of other pre-Islamic Iranian celebrations such as Now-Ruz (1st day of Spring).

\hfill

\texttt{Charlotte Cowell on 2011-12-23 at 06:11 said:}

Clearly a big day/night for the entire world and rest of the cosmos :-)

\hfill

\texttt{Perennial on 2012-12-20 at 03:12 said:}

This is from a post I made on a debate on the date of Christmas in another forum I thought I would share: The feast of Sol Invictus began the week before, and as was common in the traditional world the feast lasted several days. The day of the unconquered sun overlapped with the 25th, but did not fall on it. In addition, the feast of Sol Invictus was not widely celebrated until about 200 years after Christ's death, being popularised by Augustus Aurelian. In fact, it has been speculated that Sol Invictus was actually a rip-off of Christmas, and not the reverse. We know that the feast of Sol Invictus dates to no earlier than 274 AD at Rome. This was observed by Roman games every 4 years, evidenced by the surviving Chronology of 354, which lists the games as beginning on the 25th. However, Christmas, from ancient Martyologies, was already celebrated annually on the 25th in 336. In fact there is no reason to believe that Sol Invictus was observed December 25th prior to 354, a full 18 years after the Roman sacramentary records Christmas. In fact, scholars believe that August may have likely been the original month Sol was celebrated. Other works list solar game in October. Tertullian argued for December 25th being Christ's birth back in 198 AD, 76 years before Sol Invictus was even observed by Aurelian, so clearly none of this was new. As for the other nonsense, such as Saturnalia and the like, Saturnalia occurred in earl December, nowhere near Christmas, and ended before the 25th. Clearly, there is little connection between pagan deities and the Feast of the Nativity. The Holy Father rightly observes that Christ's birthday is dated 9 months after the Annunciation, when Our Lady conceived Our Lord. No mysteries here.

\hfill

\texttt{Cologero on 2012-12-20 at 10:29 said:}

With all due respect, Perennial, you should be seeing mysteries everywhere; this is hardly the equivalent of a birth announcement in the NY Times. Please review your rosary.

Unfortunately, you are adopting the methods of ``scientific'' history, so praised in the documents of Vatican II, but whose effects have been quite corrosive. Expresssions such as ``there is no evidence'' or ``there is no reason to believe'' are far from establishing anything concrete.

The second mistake you make is thinking Gornahoor is like any other forums you may frequent. You may want to refute Saturnalia all you want, but I don't believe that was even mentioned in the post nor in Evola's text.

I suggest you read the chapter ``The Limits of History and Geography'' from ``The Reign of Quantity'' and try to rewrite your comment in its light.

Evola's, and Guenon's, point is that dates are not arbitrary. Pace, Evola, however, Christmas is not a decomposition, but rather a deepening and a recomposition. It renewed a decomposing Tradition, which we witness in the Middle Ages. It was such, also, that a poet like Virgil could anticipate the start of a new Golden Age.

\hfill

\texttt{Perennial on 2012-12-20 at 12:56 said:}

I agree, Cologero. Obviously, this was written for a different context to be read by a different audience. I am not attempting to refute Evola, since i agree with him, and even though I do not believe that Christmas is based on the solstice, I also believe it is not arbitrary and the feast occurs where it does for a reason. I do not think it is coincidence that all of these events occur in December, either. I shared this primarily to give some factual background to parallel the post, rather than compete with it.

\hfill

\end{sffamily}
\end{footnotesize}

